% !TEX root = ../lab1_report.tex
% ==========================================================
% 实验三：发动机虚拟装配实验
% ==========================================================
\section{实验三：发动机虚拟装配实验}

\subsection{实验目的}

\begin{enumerate}
    \item 通过实验了解发动机的总体装配流程。
    \item 了解压气机、燃烧室、涡轮等大部件的安装顺序及特点。
\end{enumerate}

\subsection{实验装置}

该实验装置是一套发动机虚拟装配系统。虚拟装配系统包括讲解演示、引导装配、自主装配三种模式。参与实验的学生可首先根据讲解演示了解装配过程，之后利用引导装配熟悉装配顺序和过程，最后利用自主装配掌握整个装配流程。

\subsection{实验步骤及内容}

\subsubsection{实验准备}

实验之前要首先复习《航空发动机结构》课所学过的内容，熟悉发动机各零部件的名称、位置和功能。

\subsubsection{启动与学习}

\begin{enumerate}
    \item 启动发动机虚拟装配系统。
    \item 选择``总体装配''。
    \item 依次选择三种学习模式：
    \begin{itemize}
        \item \textbf{讲解演示}：学习发动机整体装配流程，了解各主要部件的装配顺序。
        \item \textbf{引导式训练}：按系统提示逐步完成各部件装配，熟悉装配步骤。
        \item \textbf{自主训练}：不依赖提示独立完成全部装配操作。
    \end{itemize}
\end{enumerate}

\subsubsection{装配观察要点}

在操作中注意观察各个零件的装配顺序、装配关系以及特点：
\begin{itemize}
    \item 转子组件与静子组件的对中安装方法
    \item 各级之间的定位销和螺栓连接方式
    \item 各支点轴承的安装位置、预紧力施加方式
    \item 级间密封装置（篦齿密封、石墨密封等）的安装
    \item 附件传动系统的安装接口和对齿要求
    \item 外部管路和电缆的敷设路径和固定方式
\end{itemize}

\subsection{实验分析与讨论}

\subsubsection{总体装配流程}

发动机总体装配遵循从内向外、从前向后、先冷端后热端的基本原则。本实验记录的装配流程涵盖该型发动机从核心机到外部管路，可分为以下九个阶段。

\subsubsection{第一阶段：中介机匣与传动系统装配}

安装中介机匣、内齿轮箱壳体、齿轮轴、棍棒轴承，随后依次安装低压压气机前轴承支座、高压压气机前轴承支座、高压压气机前轴。安装枢座作动盘支座及调整垫圈、枢座作动盘。安装进口导流叶片。安装各级滑油封严圈及空气封严圈。

\subsubsection{第二阶段：低压压气机装配}

安装一级低压压气机轮盘，安装低压压气机滑油管和联轴器球头。依次安装低压压气机前轴承、前空气封严圈、空气封严圈定外件、前滑油封严圈、前杯形垫圈、一级盘后隔圈、动压封严圈、一级盘前隔圈。随后按相同方式安装二至五级：前后隔圈、动压封严圈及各级轮盘。安装低压压气机后滑油封严圈、轴承座及轴承。安装低压压气机轴到低速齿轮箱传动轴的螺旋齿轮，以及后杯形垫圈。

低压静子部分：安装五级至一级低压静子叶片（由后向前）。安装止动板、低压压气机上部和下部机匣。

\subsubsection{第三阶段：前机匣与附件系统装配}

安装前轴承座、进口导流叶片、前轴承滑油泵。安装头部整流罩、进口导流叶片机匣。安装防冰空气总管、防冰空气管外罩、进气连接罩。

附件系统：安装滑油箱安装座、空气冷却滑油散热器、辅助齿轮箱、高速齿轮箱、燃油冷却滑油散热器、滑油箱。安装防喘调节器、低速齿轮箱、低压轴转速限制器。安装低压及高压测速电机、低压燃油泵。安装中介传动轴。

\subsubsection{第四阶段：高压压气机装配}

安装低压止推轴承、花键螺母、空气导管、导油环、档油套、后封严圈座、滑油分配器、前支撑管。安装滑油供油管、低压联轴器、前轴、后轴。

高压转子：安装二级高压转子盘，一至二级高压隔圈。安装一级高压静子圈、一级高压转子轮盘。二至三级静子圈和二至十二级隔圈及各级轮盘交替安装。安装高压十二级空气封严圈、大螺母。

高压静子部分：安装右部高压压气机机匣，一级至十级高压静子叶片，高压压气机机匣和静子。安装放气活门、扩散机匣、高压止推轴承座组件、出口整流叶片。

安装各级空气封严圈、滑油封严圈（前、中、后）、滑油空气管组件、滑油/通气/引气管、空气进气口。

\subsubsection{第五阶段：高压涡轮装配}

安装高压涡轮轴、高压涡轮轴承套、轴承套封严圈、空气蓖齿封严圈及封严座、后滑油蓖齿封严圈、限动环、弹性支撑。安装高压涡轮轴承、前滑油蓖齿封严环、轴承座、锁紧套筒、花键螺母、高压联轴器。安装滑油进油/回油/排气管、高压涡轮轴承支座。

\subsubsection{第六阶段：燃烧室装配}

环管燃烧室共 10 个火焰筒。先安装 1、3、5、7、9 号（奇数号）火焰筒，再安装 2、4、6、8、10 号（偶数号）火焰筒。安装燃烧室外机匣。

\subsubsection{第七阶段：涡轮导向器与转子装配}

高压涡轮：安装高压一级涡轮导向叶片、锯齿形环、扇形定位件、定位环、导向器机匣。安装高压一级涡轮叶片外圈、叶片、轮盘。安装空气封严件、外空气封严圈、封严支座、隔圈。安装高压二级导向器叶片支承座、高压二级导向叶片、涡轮盘、涡轮叶片。

低压涡轮：依次安装低压一级导向叶片、涡轮盘、涡轮叶片，以及低压二级涡轮导向叶片。安装低压涡轮轴、滑油供油管组件、低压涡轮轴承支座。

\subsubsection{第八阶段：外涵道、排气与加力系统装配}

安装引气管路、外涵道前段和后段。安装排气混合器。安装加力燃烧室和喷管。

\subsubsection{第九阶段：外部管路与控制系统装配}

安装防冰管、凸轮箱、停车开关、加力燃油调节器、压比调节器、喷口滑油泵。安装高压燃油泵、主燃油调节器、主滑油泵、低压燃油滤、分圈活门。

燃油管路连接：高压燃油泵至主燃油调节器、停车开关至副油路/主油路、主燃油调节器至停车开关、低压燃油滤至高压燃油泵。安装燃油冷却滑油散热器至低压燃油滤、至加力汽心泵的管路。安装加力燃油调节器至分圈活门、分圈活门至加力燃油总管。

最后安装发动机和喷口控制系统管路。

\subsection{结论}

通过本次虚拟装配实验，完整记录了该型发动机从核心机装配到外部管路连接的装配过程，涵盖中介机匣装配、低压/高压压气机逐级装配、环管燃烧室分奇偶号火焰筒安装、高压/低压涡轮装配、附件系统及外部管路连接等九个阶段。装配流程体现了航空发动机从内向外、先冷端后热端、逐级交替装配的基本工艺原则。三种学习模式的渐进训练使装配技能掌握更加扎实。
